\chapter*{Abstract}

\section*{\tituloPortadaEngVal}

The invention of \textbf{synthesizers} in the 1960s revolutionized the music world, previously dominated by traditional analog instruments. The arrival of \textbf{synthetic music} and the computer as a central element of modern production have defined the way artists work, modifying their sound and creative processes. Since then, music producers and sound designers have spent years learning to master synthesizers to define and create new musical timbres; however, this process often requires great mastery, as the different types of synthesis are usually complex and scalable, making it a highly complicated task for most musicians. It is common for even the best sound designers to find it difficult to \textbf{replicate other synthetic sounds}, as they depend directly on the hardware used and the sequence of effects and sound modulations applied.

That is why the application of \textbf{deep learning} to synthesis opens up a whole range of possibilities by delegating pattern detection and the technical knowledge required to \textbf{artificial intelligence}, allowing the artist to focus on the experimentation and creation that belongs to them. Through the use of \textbf{Convolutional Neural Networks (CNNs)}, this project seeks to automate parameter estimation in complex \textbf{FM synthesis}. Thus, AI solves the technical barrier of replicating a sound (\textbf{sound matching}), becoming a direct tool at the service of creativity.

\vspace{1em} 

All the developed code can be found in our GitHub repository: \\ 
\url{https://github.com/Angelji06/TFG-SoundMatchingFM}

\section*{Keywords}

\noindent Sound Matching, FM Synthesis, Deep Learning, Convolutional Neural Networks, Audio Signal Processing, Spectrograms, Python, PyTorch.



